\subsection*{NT-L15. A mean-square criterion for simple Hecke zeros and reciprocal derivatives}
Let $\eta$ be a fixed primitive unitary Hecke character of a number field $K$,
including a pure norm character. Define coefficients $b_\eta(\mathfrak a)$ by
\[
 L(s,\eta)^{-1}=\sum_{\mathfrak a\ne0}
       b_\eta(\mathfrak a)(N\mathfrak a)^{-s}
                    \qquad(\Re s>1),
 \qquad
 \mathcal M_\eta(x)=\sum_{N\mathfrak a\le x}b_\eta(\mathfrak a).
\]
Thus $b_\eta(\mathfrak a)=\mu_K(\mathfrak a)\eta(\mathfrak a)$
on ideals prime to the conductor and is zero on the other ideals.
Assume the additional weak mean-square hypothesis
\[
 B_\eta:=\limsup_{X\to\infty}\frac1{\log X}
       \int_1^X\left|\frac{\mathcal M_\eta(x)}x\right|^2dx<\infty.
                                                        \tag{NT-L15a}
\]
Then GRH holds for $L(s,\eta)$, every nontrivial zero is simple, and
\[
 \sum_{\rho=1/2+i\gamma}\frac1{|\rho L'(\rho,\eta)|^2}
                          \le B_\eta.                   \tag{NT-L15b}
\]
In particular $B_\eta>0$ and
\[
 |L'(\rho,\eta)|^{-1}\le\sqrt{B_\eta}\,|\rho|,
 \qquad |L'(\rho,\eta)|^{-1}=o(|\rho|),
                                                        \tag{NT-L15c}
\]
where the little-$o$ is along zeros of unbounded height. Also
\[
 \sum_{|\gamma|\le T}|L'(1/2+i\gamma,\eta)|^{-2}=o(T^2).
                                                        \tag{NT-L15d}
\]
No use of $A$ or $M$ is required once (NT-L15a) is assumed. Neither $A$ nor
$M$ is asserted to imply this additional mean-square hypothesis.
For $K=\mathbb Q$, $\eta=1$, the classical source is Titchmarsh,
\emph{The Theory of the Riemann Zeta-function}, 2nd ed.,
\S14.29: hypothesis (14.29.1) and Theorem14.29(A), pp.376--377;
reciprocal derivative estimate (14.29.3), p.377;
Theorem14.29(B)/(14.29.4), p.377, with proof ending p.380.
The following proof gives the Hecke formulation and its precise limsup constant.

\begin{proof}
Put $F(y)=e^{-y/2}\mathcal M_\eta(e^y)$ for $y\ge0$ and
$A(Y)=\int_0^Y|F(y)|^2dy$. Changing variables $x=e^y$ gives
$B_\eta=\limsup_{Y\to\infty}A(Y)/Y$.
For every $\epsilon>0$ there is $C_\epsilon$ such that
\[
 A(Y)\le (B_\eta+\epsilon)Y+C_\epsilon\qquad(Y\ge0).
                                                        \tag{NT-L15e}
\]
In particular the Laplace integral
\[
 H(w)=\int_0^\infty F(y)e^{-wy}\,dy
\]
is holomorphic for $\Re w>0$. To check this, integrate $|F|^2$ against
$e^{-\sigma y}$ by parts using (NT-L15e), then apply Cauchy--Schwarz to
$F(y)e^{-wy}$. The same argument with any fixed factor $y^j$ proves
uniform convergence of every derivative on each compact subset of that
half-plane. For $\Re w>1/2$, absolute convergence of the inverse Euler
series and partial summation give
\[
 H(w)=\int_1^\infty\mathcal M_\eta(x)x^{-(w+1/2)-1}\,dx
      =\frac1{(w+1/2)L(w+1/2,\eta)}.                    \tag{NT-L15f}
\]
The meromorphic identity theorem extends (NT-L15f) throughout $\Re w>0$.
Since $H$ is holomorphic there and $w+1/2\ne0$, $L(s,\eta)$ has no zeros
with $\Re s>1/2$. The conjugate character has coefficients
$b_{\bar\eta}(\mathfrak a)=\overline{b_\eta(\mathfrak a)}$, hence the
same mean-square hypothesis and the same zero-free half-plane.
The Hecke functional equation now excludes nontrivial zeros of
$L(s,\eta)$ with $\Re s<1/2$. This proves GRH, including non-real
characters; possible pure-norm poles only give zeros of $H$ and cause
no difficulty.

For $\sigma>0$ introduce the positive inner product
\[
 \langle f,g\rangle_\sigma
       =\sigma\int_0^\infty f(y)\overline{g(y)}e^{-\sigma y}dy.
\]
An integration by parts and (NT-L15e) show that
\[
 \|F\|_\sigma^2
   =\sigma^2\int_0^\infty A(y)e^{-\sigma y}dy,
 \qquad \limsup_{\sigma\downarrow0}\|F\|_\sigma^2\le B_\eta.
                                                        \tag{NT-L15g}
\]
Each $e_t(y)=e^{ity}$ has norm one. Consequently
\[
 |\sigma H(\sigma+it)|
       =|\langle F,e_t\rangle_\sigma|
       \le\|F\|_\sigma.
                                                        \tag{NT-L15h}
\]
At a zero $\rho=1/2+i\gamma$, (NT-L15f) has a pole of the same order
as that zero, since $\rho\ne0$. The bound (NT-L15h) as
$\sigma\downarrow0$ rules out order greater than one. Thus every zero
is simple and
\[
 \lim_{\sigma\downarrow0}\langle F,e_\gamma\rangle_\sigma
       =a_\gamma:=\frac1{\rho L'(\rho,\eta)}.
                                                        \tag{NT-L15i}
\]
For two distinct ordinates,
\[
 \langle e_\gamma,e_\nu\rangle_\sigma
       =\frac\sigma{\sigma-i(\gamma-\nu)}\longrightarrow0;
\]
the same inner product is one when $\gamma=\nu$.
Take an arbitrary finite set $S$ of zero ordinates and put
$P_S=\sum_{\gamma\in S}a_\gamma e_\gamma$.
Expanding the nonnegative number $\|F-P_S\|_\sigma^2$, using
(NT-L15g)--(NT-L15i), and taking the upper limit gives
\[
 0\le\limsup_{\sigma\downarrow0}\|F-P_S\|_\sigma^2
       \le B_\eta-\sum_{\gamma\in S}|a_\gamma|^2.
\]
The supremum over finite $S$ proves (NT-L15b). Nontrivial Hecke zeros
exist, by the classical zero-count asymptotic, so $B_\eta>0$.
Every summand is at most $B_\eta$, proving the first assertion of
(NT-L15c). Square summability makes $a_\gamma\to0$ along unbounded
ordinates and proves the second assertion. Finally, for $T\ge1$,
\[
 T^{-2}\sum_{|\gamma|\le T}|L'(\rho,\eta)|^{-2}
   =\sum_\gamma {\bf1}_{|\gamma|\le T}
           \frac{1/4+\gamma^2}{T^2}|a_\gamma|^2.
\]
Each factor before $|a_\gamma|^2$ tends to zero and is at most $5/4$.
Dominated convergence for the summable series proves (NT-L15d).
\end{proof}

The finite Bessel argument proves an inequality. It does not assert an
exact mean-square limit, completeness of an infinite Fourier expansion
of $F$, or a limiting distribution for $F$. Such assertions require
additional approximation estimates.
Geometrically, the coefficients are the inverse Euler weights of the
rank-one arithmetic object attached to $\eta$; squarefree ideals are
squarefree effective divisors on $\operatorname{Spec}\mathcal O_K$.
The new hypothesis is analytic control of these weighted divisors and
is not a consequence of a finite-field purity theorem established here.

\subsection*{NT-Z19. Weak mean square supplies the canonical arithmetic basis gate}
Use the canonical completed, pole-removed $\Xi_\eta$, inverse Mellin
function $\psi_\eta$, and closed quotient
\[
 Q_\eta=\mathcal E/\overline{\psi_\eta*\mathcal E},\qquad
 \mathcal E=\{f\in C^\infty(\mathbb R):
      \sup_x e^{N|x|}|f^{(j)}(x)|<\infty\text{ for all }N,j\ge0\}
\]
of MT-M12/MT-M16. Under (NT-L15a), all zeros are simple and
\[
 \overline{\psi_\eta*\mathcal E}
   =\{f:\mathcal Mf(\rho)=0\text{ for every nontrivial zero }\rho\},
 \qquad Q_\eta\simeq s(Z).
                                                        \tag{NT-Z19a}
\]
Here $\mathcal Mf(s)=\int_\mathbb Rf(x)e^{sx}dx$,
and $s(Z)$ is the space of sequences decreasing faster than every
power of $1+|\operatorname{Im}\rho|$. The quotient has an absolute
Schauder basis with these Mellin evaluations as coefficient functionals.
Convolution by every $h\in\mathcal E$ is nuclear. In the diagonal
Hilbert realization $\ell^2(Z)$ it is trace class, with Hilbert trace
\[
 \operatorname{tr}_{\ell^2(Z)}T_h=\sum_\rho\mathcal Mh(\rho).
\]
Writing $h^\#(x)=e^{-x}\overline{h(-x)}$, the operator
$T_{h*h^\#}$ is positive and
\[
 \operatorname{tr}_{\ell^2(Z)}T_{h*h^\#}
                     =\sum_\rho|\mathcal Mh(\rho)|^2\ge0.
                                                        \tag{NT-Z19b}
\]
\begin{proof}
NT-L15 proves GRH, simplicity, and
$|L'(\rho,\eta)|\ge B_\eta^{-1/2}|\rho|^{-1}$.
Since $|\rho|\le1+|\operatorname{Im}\rho|$, this is precisely the
polynomial derivative lower bound of MT-M12, or condition (1) of
MT-M16 with every multiplicity equal to one. The entire proofs of
those results give (NT-Z19a) and the absolute basis. MT-M13/MT-M17
then give nuclearity and the trace-class Hilbert realization. For
completeness, on each coordinate the convolution eigenvalue is
$\mathcal Mh(\rho)$, and direct substitution gives
$\mathcal M h^\#(s)=\overline{\mathcal Mh(1-\bar s)}$.
GRH implies $1-\bar\rho=\rho$, so the product eigenvalue is
$|\mathcal Mh(\rho)|^2$. Rapid Mellin decay and NT-Z10's local zero
count make both displayed trace series absolutely convergent. This
proves positivity and (NT-Z19b).
\end{proof}
This is a conditional analytic realization of the zero divisor with an
explicit arithmetic gate. It does not identify the quotient with an
independently constructed motivic or étale cohomology group.
For the Riemann zeta function the sole additional arithmetic condition
is the classical weak Mertens conjecture; $A$ is not needed separately.
